\documentclass[12pt]{article}
\usepackage{enumitem}
\usepackage{geometry}
\usepackage{amsmath, amssymb}
\geometry{margin=1in}
\begin{document}
\begin{center}
Astronomy - Graduate Level Assessment 1 \
Total Marks: 40 \
Answer all questions.
\end{center}

\section*{Multiple Choice (10 marks)}
\begin{enumerate}[label=\arabic*.]
\item According to the Solar Nebular theory what are asteroids and comets?
\begin{enumerate}[label=(\alph*)]
    \item They are the shattered remains of collisions between planets.
    \item They are chunks of rock or ice that condensed long after the planets and moons had formed.
    \item They are chunks of rock or ice that were expelled from planets by volcanoes.
    \item They are leftover planetesimals that never accreted into planets.
\end{enumerate}
\item The so-called "bigfoot" on Mars was actually a rock that was about 5 cm tall. It had an angular size of about 0.5 degrees (\textasciitilde{}30 pixels). How far away was this...
\begin{enumerate}[label=(\alph*)]
    \item About 6 meters
    \item About 6 feet
    \item About 10 meters
    \item About 10 feet
\end{enumerate}
\item Which of the following is/are true?
\begin{enumerate}[label=(\alph*)]
    \item Titan is the only outer solar system moon with a thick atmosphere
    \item Titan is the only outer solar system moon with evidence for recent geologic activity
    \item Titan's atmosphere is composed mostly of hydrocarbons
    \item A and D
\end{enumerate}
\item Which of the following most likely explains why Venus does not have a strong magnetic field?
\begin{enumerate}[label=(\alph*)]
    \item Its rotation is too slow.
    \item It has too thick an atmosphere.
    \item It is too large.
    \item It does not have a metallic core.
\end{enumerate}
\item Which of the following was not cited as evidence for life in the martian meteorite ALH84001?
\begin{enumerate}[label=(\alph*)]
    \item Complex organic molecules specifically PAHs
    \item Magnetite grains similar to those formed by bacteria on Earth
    \item Carbonate minerals indicating a thicker warmer Martian atmosphere
    \item Amino acids with a preferred orientation or "chirality"
\end{enumerate}
\end{enumerate}

\section*{True / False (10 marks)}
\textit{Answer True or False and justify briefly.}
\begin{enumerate}[label=\arabic*.]
\item True or False: The correct answer to 'The axis of the Earth is tilted at an angle of approximately ... relative to the orbital plane around the Sun.' is '20.3 degrees'.
\item True or False: The correct answer to 'Which of the following is/are common feature(s) of all fresh (i.e. not eroded) impact craters formed on solid surfaces:' is 'A and B only'.
\item True or False: The correct answer to 'Most rocks on the Moon's surface are older than those on the Earth's surface. The best evidence for this is:' is 'The Moon's surface is more heavily eroded than the Earth's surface.'.
\item True or False: The correct answer to 'The constellation ... is a bright W-shaped constellation in the northern sky.' is 'Cassiopeia'.
\item True or False: The correct answer to 'One astronomical unit (AU) is equal to approximately ...' is '150 million km'.
\end{enumerate}

\section*{Long Form Response (20 marks)}
\textit{Provide thorough reasoning and reference key steps.}
\begin{enumerate}[label=\arabic*.]
\item Why do we look for water-ice in craters at Mercury's pole?
\item By locating the north celestial pole (NCP) in the sky how can you determine your latitude?
\end{enumerate}

\vfill
\noindent\textit{End of Paper}
\end{document}